\documentclass[10pt,a4paper]{amsart}
\usepackage[margin=19mm]{geometry}
\usepackage{amsmath,amssymb,mathtools}
\usepackage[scr=rsfs,cal=euler]{mathalfa}
\usepackage{xfrac}
\usepackage[unicode,hidelinks,bookmarks=false]{hyperref}
\newcommand{\ie}{i.e.\ }
\newcommand{\cf}{cf.\ }
\newcommand{\resp}{respectively\ }
\newcommand{\Subsection}{\subsection}
\providecommand{\nobreakdash}{\nobreak\hbox{-}}
\setlength{\emergencystretch}{2em}
\multlinegap0pt
\usepackage{bm}
\usepackage{xcolor}
\usepackage{amssymb}
\usepackage{float}
\newtheorem{lemma}{Lemma}
\newtheorem{corollary*}{Corollary}
\title{Almost Bayesian: Dynamics of SGD Through Singular Learning Theory}
\author{Max Hennick, Stijn De Baerdemacker}
\date{Source: arXiv:2503.22478v2 (CC BY 4.0)}
\begin{document}
\maketitle

\noindent\emph{Source and licence.} Almost Bayesian: Dynamics of SGD Through Singular Learning Theory. Version arXiv:2503.22478v2 (CC BY 4.0), distributed by arXiv under the Creative Commons Attribution 4.0 International licence, http://creativecommons.org/licenses/by/4.0/, as recorded in the arXiv licence metadata for this exact version and shown on the abstract page https://arxiv.org/abs/2503.22478v2. Full licence notice: \texttt{licenses/arxiv-2503.22478-CC-BY-4.0.txt}.\par\smallskip

\noindent\textit{Editorial scope.} Counted unit: the article's corollary* corollary* (source0.tex lines 1313--1318). The article's complete original proof of it is delivered with it (source0.tex lines 1319--1345, one \texttt{proof} environment of 27 lines). No mathematics is written, repaired or re-derived: the statement and the proof are the article's own lines, in the article's own notation, and no retained line is substituted or re-flowed.  Retained uncounted as the source-local context the proof needs: lines 846--848.  Nothing was omitted from any retained span, so the package carries no omission record.  No retained line contains a citation, so the wrapper carries no bibliography and no bibliography entry.  No retained line contains a figure environment, an image-inclusion command or drawing code, so no image is delivered and the package declares no asset dependency.  Editorial formatting only, in addition: an AMS \texttt{amsart} preamble that restates the article's own declarations and macros used by the retained lines, namely the environments \texttt{lemma}, \texttt{corollary*} on separate counters, together with the standard packages those lines need.  The retained environments print in this document as Lemma 1, Corollary* 1 in order of appearance, whereas the article prints the same items with the numbering of its own full document.  The complete native source of the article as distributed in the arXiv e-print for this exact version is bundled unchanged as \texttt{sources/175558/source0.tex}.
\begin{lemma}\label{lem:spect_ineq}
Suppose the loss function $\mathcal{L}$ is non-convex and non-constant on $W$. Then with spectral dimension $d_s$ as $t \to \infty$ with fractal dimension $\lambda(w(t))$ on $\mathcal{W} \subset W$, the inequality $d_s \leq \lambda(w(t))$ holds (in the small learning rate regime).
\end{lemma}

\begin{corollary*}
For time $t$ as $t \to \infty$, we have $d_s \leq \bar{\lambda}(w(t))$ where
    \begin{equation}
        \bar{\lambda}(w(t)) = \lim_{\tau \to \infty} \frac{1}{\tau} \int_0^\tau \lambda(w(t)) dt
    \end{equation}
\end{corollary*}

\begin{proof}
Let $\tau_0$ be the time such that for all $\tau > \tau_0$, the inequality of lemma \ref{lem:spect_ineq} holds. Consider then a time $T \gg \tau_0$ and consider the integral
\begin{equation}
\int_0^T \lambda(w(t)) dt = \int_0^{\tau_0} \lambda(w(t)) dt + \int_{\tau_0}^T \lambda(w(t)) dt
\end{equation}
and since $\tau_0$ is finite we can take the first portion of this integral to be a constant (since we know that the LLC is bounded above by $\frac{d}{2}$ where $d$ is the number of free parameters):
\begin{equation}
    \int_0^{\tau_0} \lambda(w(t)) dt = C
\end{equation}
By the result of lemma \ref{lem:spect_ineq} we have that for all times greater than $\tau_0$, we must have
\begin{equation}
    \int_{\tau_0}^T \lambda(w(t)) dt \geq \int_{\tau_0}^T d_s dt
\end{equation}
and since $d_s$ is constant
\begin{equation}
    \int_{\tau_0}^T \lambda(w(t)) dt \geq (T-\tau_0)d_s
\end{equation}
which means that by adding $C$ to both sides and dividing by $T$ we get
\begin{equation}
    \frac{1}{T}\int_0^T \lambda(w(t)) dt \geq \frac{(T-\tau_0)d_s}{T} + \frac{C}{T}
\end{equation}
From this we get
\begin{equation}
    \frac{1}{T}\int_0^T \lambda(w(t)) dt \geq d_s + \frac{(C-\tau_0)d_s}{T}
\end{equation}
where the term $\frac{(C-\tau_0)d_s}{T}$ vanishes as $T \to \infty$ since $C$ must be finite.
\end{proof}

\end{document}
